\section{Quantum Reservoirs and Reset-Window Models}
\label{sec:models}

A QRC processes inputs one at a time. At discrete step $t$, input $a_t$
updates the previous state, and selected measurements provide the output
\cite{fujii2017harnessing,mujal2021opportunities}:
\[
 \rho_t=\Phi_{a_t}(\rho_{t-1}),
 \qquad y_t=\mathcal C(\rho_t).
\]
Here $\rho_t$ is the carried quantum state; the quantum channel $\Phi_{a_t}$
describes its physical evolution, including dissipation.
The fixed map $\mathcal C$ collects $m$ observable expectations
$\operatorname{Tr}(O_j\rho_t)$, $j=1,\ldots,m$: average measurement values,
shown as angle brackets in Figure~\ref{fig:information-interface}.
A trained linear readout can be included in $\mathcal C$ by absorbing its
fixed weights into the observables.

A reset-window QELM starts from the same quantum state for every prediction
and encodes the $d$ most recent inputs
\cite{innocenti2023potential,kawanabe2026tdqelm}.
For a history $h$ ending at $t$, with target output $y(h)=y_t$, its window is
\[
 S_d(h)=(a_{t-d+1},\ldots,a_t).
\]
The encoder $E(S_d(h))$ produces quantum features: observable expectations
of $P_j$, $j=1,\ldots,B$, including measurement checkpoints.
Thus $B$ counts measured expectations, not qubits.
A trained linear readout combines them. With exact expectations, the
complete prediction is
\[
 \widehat y_t=f(S_d(h)),
\]
where $f$ includes the quantum feature map and readout.

Identical windows therefore produce identical predictions for a
fixed model, even if earlier inputs differ.
With finitely sampled measurements, predictions can fluctuate, but their
probability distributions remain identical.
More qubits or features enrich the window function, but cannot recover
an earlier input that the window leaves unspecified.

For emulation we hold $\mathcal C$ fixed and target $y_t$, not the full state.
The emulation simulations use observable expectations directly; the separate
forecasting tests fit both readouts to external task targets.

Figure~\ref{fig:information-interface}'s output plots are schematic, not
fitted QELM predictions. The coral prediction is the conditional mean over two equally
likely old-input values with the same later inputs and known baseline.
Above, the old input is fully forgotten by window expiry, a special case
of fading memory. Below, later inputs gradually reveal retained information.
Sufficiency here concerns these example histories, not uniform accuracy
as the window length grows; Section~\ref{sec:criterion}
gives that criterion and the exact scalar error bound $D_d/2$. In the input
plot, $k$ indexes displayed steps and $t$ marks the current one.
